% Lösungen Einheit 1 von 3 – Kreisumfang (zusammenbau v0.1)
\einheitenkopf{Einheit 1 von 3 · Kreisumfang}

\erg{9}{a) Umfang}
\erg{10}{a) Durchmesser; $r = 4{,}5$ cm \quad b) $u = \pi \cdot 6 \approx 18{,}8$ cm \quad c) $u = 2 \cdot \pi \cdot 4 \approx 25{,}1$ cm \quad d) $u = \pi \cdot 4{,}6 \approx 14{,}45$ cm \quad e) $d = 47 : \pi \approx 15{,}0$ cm \quad f) $d = 75 : \pi \approx 23{,}87$ cm, $r \approx 11{,}9$ cm \quad g) $\pi \cdot 12 : 2 \approx 18{,}8$ m \quad h) $u = \pi \cdot 66 \approx 207{,}3$ cm; Weg $= 250 \cdot u \approx 51\,836$ cm $\approx 518{,}4$ m \quad i) Länge $= 2 \cdot \pi \cdot 3{,}6 \approx 22{,}6$ cm, Breite $= 11$ cm}
\erg{11}{a) Strecke von $P$ durch $M$ bis zum gegenüberliegenden Rand, beschriftet mit $d$}

11a)

\begin{kreis}[2] \mittelpunkt{M} \kreispunkt{30}{P} \durchmesser{30}{d} \end{kreis}

\erg{12}{a) Zirkel auf $r = 3$ cm einstellen, Kreis um $M$ zeichnen}

12a)

\begin{kreis}[3] \mittelpunkt{M} \radius{30}{$r$} \end{kreis}

\erg{13}{a) $22 : 7 \approx 3{,}14$}
\erg{14}{a) Ben hat den Durchmesser als Radius eingesetzt. Richtig: $u = \pi \cdot 18 \approx 56{,}5$ cm \quad b) Alle Kreise haben dieselbe Form: Wird der Durchmesser doppelt oder dreimal so groß, wird der Umfang ebenso doppelt oder dreimal so groß – der Quotient bleibt gleich, nämlich $\pi$. \quad c) $d = 283 : \pi \approx 90{,}1$ cm; das ist mehr als $75$ cm, die Eiche steht unter Schutz. \quad d) $3{,}1$; $6{,}3$; $9{,}4$; $18{,}8$}
